% Generated draft: requires mathematical and source review.
\documentclass[11pt]{article}
\usepackage[margin=25mm]{geometry}
\usepackage{fontspec}
\setmainfont{Latin Modern Roman}
\usepackage{amsmath,amssymb,mathtools,unicode-math}
\setmathfont{Latin Modern Math}
\usepackage{enumitem,needspace,xurl,hyperref}
\setlength{\parindent}{0pt}
\setlength{\parskip}{5pt}
\setlength{\emergencystretch}{3em}
\newcommand{\N}{\mathbb{N}}
\newcommand{\Z}{\mathbb{Z}}
\newcommand{\Q}{\mathbb{Q}}
\newcommand{\R}{\mathbb{R}}
\newcommand{\C}{\mathbb{C}}
\newcommand{\F}{\mathbb{F}}
\newcommand{\A}{\mathbb{A}}
\newcommand{\PP}{\mathbb{P}}
\newcommand{\E}{\mathbb{E}}
\newcommand{\eps}{\varepsilon}
\newcommand{\dd}{\,\mathrm{d}}
\begin{document}
\section*{UnsolvedMath ID: 80}
\textbf{Source problem code:} catalogue entry
\subsection{Definitions and mathematical statement}
The source attributes ``Optimal Nanostructures'' to DARPA's Mathematical Challenge Eleven (2007). Its objective is to develop mathematics for creating optimal symmetric structures by nanoscale self-assembly. The supplied review describes the original challenge as a globally symmetric local-rule formulation; the background places it at the intersection of crystallography, optimization, and molecular dynamics.

In this context, local rules describe how constituent building blocks interact or attach, while global symmetry concerns the assembled structure as a whole. For a specified configuration space $X$ and a specified group $G$ acting on $X$, invariance of a configuration $x$ under that group means
\[
 g\cdot x=x\qquad\text{for every }g\in G.
\]
This is only a definition of a possible symmetry requirement: the source does not specify $X$, $G$, whether the target has all or only some of the proposed symmetries, or even a particular geometric or combinatorial model of an assembled structure.

\textbf{Source-faithful research objective.} Develop mathematical principles and design methods for local self-assembly rules that create globally symmetric nanoscale structures and achieve an appropriate optimality criterion. The source gives no objective function or ordering by which an assembly is optimal. It also supplies no building-block model, admissible rule class, assembly dynamics, constraints, or treatment of errors. Thus no optimization problem, existence theorem, or guarantee of successful assembly can be uniquely quantified from this record.

A precise version requires the self-assembly formalism, target symmetry, admissible designs and structures, optimization objective, and error or success criterion to be fixed. These are unresolved inputs, not hypotheses silently imposed in this draft. In particular, the record does not select deterministic versus stochastic assembly or require an almost-sure, high-probability, or worst-case success guarantee.

The literature triage dated 2026-08-17 calls this a collection of optimization and physical-design problems, states that nanoscale self-assembly and symmetric-structure design remain active research areas, and says a full solution or refutation is inapplicable without a model of local rules, objective, and assembly errors. Its \emph{partially solved} classification and the registry's \emph{open} label are supplied descriptions of a research agenda, not a verified status for a single proposition.
\subsection{Short English statement}
Find mathematical principles for designing local interactions that assemble nanoscale building blocks into globally symmetric structures, optimally according to a specified criterion. The source leaves the assembly model, symmetry, objective, and handling of errors unspecified, so it defines a research agenda rather than one theorem.
\subsection{Sources}
\begin{itemize}
\item UnsolvedMath local registry, data/unsolvedmath/problems.json, permanent ID 80, source code catalogue entry, ``Optimal Nanostructures,'' attributed to DARPA (2007).
\item DARPA DSO, Broad Agency Announcement 07-68, Mathematical Challenge Eleven (2007), supplied reference: \url{https://web.math.utk.edu/~vasili/refs/darpa07.MathChallenges.html}. The supplied literature triage identifies it as a primary source and cites the globally symmetric local-rule formulation. The link has not been independently consulted for this draft.
\item Ulam AI and the UnsolvedMath Contributors, \emph{UnsolvedMath: A Curated Collection of Open Mathematics Problems}, including this local record's literature triage dated 2026-08-17. Dataset curation and original metadata are supplied under CC BY 4.0; referenced source material retains its own terms.
\end{itemize}
\end{document}
